% Source PDF: source/普通高考/2012/2012山东文.pdf, page 1 (question 7).
% PDF embedded-bitmap attempt: pdfimages -list found no image objects; the flowchart is vector line art.
% Grid evidence: tmp/redraw_flowchart_2012_shandong_liberal/source/page-1-grid100.png (100 px major grid).
% No-grid comparison crop: tmp/redraw_flowchart_2012_shandong_liberal/source/q7_original_tight.png.
% Elements: 开始; 输入 a; P=0,Q=1,n=0; decision P≤Q; updates P=P+a^n, Q=2Q+1, n=n+1;
% 输出 n; 结束; branch labels 是/否.
% Node -> directed-edge list: 开始→输入a→P=0,Q=1,n=0→P≤Q;
% 是→P=P+a^n→Q=2Q+1→n=n+1→left return→P≤Q;
% 否→输出n→结束. The loop return enters the decision stem above the diamond;
% the output branch is independent and shares no line with the loop.
% solid segments: all node borders and directed connectors; dashed segments: none.
% labels: every node is locally zero-indented and centered; 是/否 sit beside their exits.

\begin{tikzpicture}[
  >=Stealth,
  line width=0.50pt,
  line cap=round,
  line join=round,
  font=\normalsize,
  flowbox/.style={draw, minimum height=0.68cm, inner xsep=4pt, inner ysep=2pt,
    anchor=center, align=center, execute at begin node={\setlength{\parindent}{0pt}}},
  terminal/.style={flowbox, rounded corners=0.18cm, minimum width=1.30cm,
    minimum height=0.72cm},
  io/.style={flowbox, trapezium, trapezium left angle=72, trapezium right angle=108,
    minimum height=0.78cm},
  decision/.style={flowbox, diamond, aspect=2.8, inner sep=2pt,
    minimum width=3.45cm, minimum height=0.94cm},
]
  \node[terminal] (start) at (0,9.35) {开始};
  \node[io, minimum width=1.75cm] (input) at (0,8.15) {输入 $a$};
  \node[flowbox, minimum width=3.55cm] (init) at (0,6.92)
    {$P=0,Q=1,n=0$};
  \coordinate (merge) at (0,6.28);
  \node[decision] (test) at (0,5.20) {$P\leq Q$};

  \node[flowbox, minimum width=2.75cm] (pupdate) at (0,3.82)
    {$P=P+a^n$};
  \node[flowbox, minimum width=2.60cm] (qupdate) at (0,2.50)
    {$Q=2Q+1$};
  \node[flowbox, minimum width=2.35cm] (nupdate) at (0,1.18)
    {$n=n+1$};

  \node[io, minimum width=1.75cm] (output) at (3.55,3.82) {输出 $n$};
  \node[terminal] (finish) at (3.55,2.50) {结束};

  \draw[->] (start.south) -- (input.north);
  \draw[->] (input.south) -- (init.north);
  \draw (init.south) -- (merge);
  \draw[->] (merge) -- (test.north);

  % 是 branch: update chain followed by an independent left return loop.
  \draw[->] (test.south) -- node[right=2pt] {是} (pupdate.north);
  \draw[->] (pupdate.south) -- (qupdate.north);
  \draw[->] (qupdate.south) -- (nupdate.north);
  % Put the return arrowhead on the outer horizontal leg; the merge itself is
  % an unarrowed degree-3 connector into the decision stem.
  \draw (nupdate.west) -- (-3.35,1.18);
  \draw (-3.35,1.18) -- (-3.35,6.28);
  \draw[->] (-3.35,6.28) -- (-0.60,6.28);
  \draw (-0.60,6.28) -- (merge);

  % 否 branch: output and finish are separate from the loop path.
  \draw[->] (test.east) -- node[above] {否} (3.55,5.20) -- (3.55,4.28) -- (output.north);
  \draw[->] (output.south) -- (finish.north);
  \path[use as bounding box] (-3.65,0.70) rectangle (4.65,9.72);
\end{tikzpicture}
